\bookchapter{제46장: 심화 여정——사용자에서 아키텍트로}

\chapterintro{37장부터 44장까지의 방법론을 익혀 AI로 프로젝트 개발을 능숙하게 해 낼 수 있게 됐다면, 다음 단계는 무엇일까요? 이번 장에서는 'AI 도구 사용자'에서 'AI 시스템 아키텍트'로 나아가는 심화 여정을 살펴봐요. 이 단계의 관심사는 더 이상 'AI에게 어떻게 좋은 코드를 짜게 만들까'가 아니라 '개발 전체 흐름에 AI를 어떻게 더 잘 녹여낼까'예요.}

\booksection{MCP 프로토콜과 도구 확장}

MCP(Model Context Protocol)는 2025년에 등장해 2026년에 점점 널리 퍼진 개방형 프로토콜로, AI 모델과 외부 도구 사이의 통신 표준을 정의해요. 쉽게 말해 MCP는 AI가 '도구를 쓸 수 있게' 해 줘요——코드를 생성할 뿐 아니라 명령을 실행하고, 데이터베이스를 조회하고, API를 호출하고, 파일 시스템을 다룰 수도 있죠.

MCP 이전에는 AI 도구마다 도구 통합 방식이 달라서 서로 호환되지 않았어요. MCP의 등장이 뜻하는 건 이거예요: MCP 표준에 맞는 도구 하나면 MCP를 지원하는 어떤 AI 도구에서도 쓸 수 있다는 거죠. 덕분에 도구 통합 비용이 크게 낮아졌고, AI의 능력 범위도 '텍스트 생성'에서 '동작 실행'으로 넓어졌어요.

심화 사용자에게 MCP의 의미는 AI를 위한 맞춤형 도구 능력을 만들 수 있다는 데 있어요. 예를 들어 팀에 자체 배포 시스템이 있다면 MCP 도구를 하나 작성해서 AI가 배포 시스템을 직접 호출하게 할 수 있고, 사람이 일일이 터미널로 전환해서 작업할 필요가 없어지죠. 이런 'AI가 도구를 호출하는' 능력이야말로 Agentic Engineering의 기반 시설이에요.

\booksection{커스텀 워크플로와 Skills}

정해진 AI 코딩 절차를 반복해서 수행하게 된다면, 그 절차를 '커스텀 워크플로'나 'Skills'로 굳혀 두면 돼요. 예를 들어 팀에 먼저 보안을 확인하고, 다음으로 성능을 확인하고, 이어서 유지보수성을 확인하는 표준 코드 리뷰 절차가 있다면, 이 절차를 Skill 하나로 만들어서 AI가 코드 리뷰를 할 때마다 자동으로 그 순서대로 수행하게 할 수 있어요.

커스텀 워크플로의 핵심 가치는 '일관성'이에요——실행할 때마다 같은 기준을 따르도록 보장해 주니까, 대화가 길거나 맥락이 달라도 어긋나지 않아요. 팀 협업에서는 특히 중요해요: 모든 멤버가 같은 워크플로를 쓰면 나오는 코드의 스타일과 품질 기준도 한결같아지거든요.

커스텀 워크플로를 만드는 방법은 도구마다 달라요. Cursor는 Rules와 Custom Instructions를 지원하고, Claude Code는 커스텀 명령을, Replit은 Runme를 지원하죠. 하지만 핵심 발상은 같아요: 여러분의 모범 사례를 AI가 실행할 수 있는 지시 템플릿으로 '코드화'하는 거예요.

\booksection{멀티 에이전트 협업 방식}

2026년에 들어서며 멀티 에이전트 협업이 뜨는 화두가 됐어요. 핵심 발상은 작업을 한 AI에게만 맡기지 않고, 역할을 나눈 여러 AI가 협력하게 한다는 거예요. 예컨대 한 AI는 코드 작성, 한 AI는 테스트 작성, 한 AI는 코드 리뷰를 맡는 식으로——각자 제 소임을 다하면서 서로를 검사해요.

멀티 에이전트 협업의 장점은 '분업'과 '교차 검증'이에요: 각 에이전트가 제 분야에 집중하니 산출물 품질이 높아지고, 에이전트 사이의 교차 검사로 혼자라면 놓쳤을 문제를 발견할 수 있죠. 하지만 단점도 뚜렷해요: 조율 비용이 크고, 통신 오버헤드가 크고, 디버깅이 더 복잡해요.

\begin{tcbinfo}{작은 팁: 혼자서도 '멀티 에이전트'를 즐길 수 있어요}
개인 개발자라면 복잡한 멀티 에이전트 시스템을 구축할 필요 없어요——여러 역할을 '연기'하면서 멀티 에이전트 협업을 흉내 내면 돼요. 예를 들어: 먼저 한 대화에서 AI에게 기능 코드를 작성하게 하고; 이어서 새 대화를 열어(CLAUDE.md 맥락을 제공하고) AI에게 '코드 리뷰어로서 이 코드를 검토해 줘'라고 하고; 마지막으로 또 하나의 대화를 열어 '리뷰 의견에 따라 코드를 고쳐 줘'라고 하면 돼요. 이런 '혼자 여러 역할' 방식은 복잡한 기반 시설 없이도 멀티 에이전트 협업의 이점 대부분을 누리게 해 줘요.

\end{tcbinfo}

\booksection{꾸준히 배우는 길}

AI 코딩은 빠르게 변하는 분야예요. 오늘의 모범 사례가 내일이면 갱신이 필요할지도 몰라요. 꾸준히 배우기 위한 제안 몇 가지예요:

첫째, 가공된 해석이 아니라 원 출처를 보세요. Andrej Karpathy의 원본 트윗과 CLAUDE.md, Anthropic과 OpenAI의 공식 문서, 주요 도구의 업데이트 로그——이것들이 1차 정보예요. 커뮤니티 글과 튜토리얼은 이런 정보를 해석하고 요약한 것이라 가치가 있지만 왜곡될 수도 있어요. 원본 자료를 읽는 습관을 들이세요.

둘째, 실천으로 배우기예요. '다 배우고 나서 하기'가 아니라 '하면서 배우세요'. 새로운 문제 유형을 만날 때마다(예: 파일 업로드를 처음 다룰 때, 외부 API를 처음 붙일 때) 그것을 기회 삼아 관련 지식을 배우는 거예요. 실전에서 배우는 효율이 상황과 동떨어진 이론 공부보다 훨씬 높아요.

셋째, 도구에 대한 감각은 유지하되 새것만 좇지는 마세요. 한두 달에 한 번은 반나절을 들여 최신 도구와 트렌드를 살펴보되, 새 도구가 나올 때마다 갈아타지는 않는 거예요. 판단 기준은 이것이에요: 새 도구가 여러분의 어떤 병목을 뚜렷하게 개선해 주는가? 그렇지 않다면 지금 조합으로 충분해요.

넷째, 개인 지식 창고를 만드세요. 복잡한 문제를 해결했던 대화, 함정에 빠졌을 때의 교훈, 효과가 있었던 방법을 기록하세요. 공개해서 알릴 필요는 없어요——이건 여러분만의 '방법론 노트'예요. 시간이 쌓이면 가장 소중한 자산이 될 거예요.

\begin{tcbpractice}{따라 해 봐요: '실수 노트' 만들기}
오늘 문서를 하나 만드세요(메모 앱이든 .md 파일이든 좋아요). 그리고 지금까지 배운 AI 코딩 교훈 중 가장 기억하고 싶은 세 가지를 적으세요——'묻기 전에 먼저 생각 정리하기' 같은 것도 좋아요. 앞으로 함정에 빠질 때마다, 새로운 노하우를 배울 때마다 한 줄씩 기록하세요. 반년 뒤에 돌아보면 오늘의 자신에게 고마워할 거예요.

\end{tcbpractice}

\begin{tcbsummary}{이번 장 정리}
\begin{itemize}
  \item MCP는 AI와 도구 사이의 통신 표준: AI가 명령 실행, 데이터베이스 조회, API 호출을 할 수 있게 해 줘요
  \item 커스텀 워크플로/Skills의 핵심 가치는 '일관성': 모범 사례를 지시 템플릿으로 코드화해요
  \item 멀티 에이전트 = 분업 + 교차 검증. 개인 개발자는 '혼자 여러 역할'로 흉내 낼 수 있어요
  \item 꾸준히 배우는 네 가지 방법: 원 출처 읽기, 하면서 배우기, 새것만 좇지 않기, 개인 지식 창고 만들기
\end{itemize}

\end{tcbsummary}
